\chapter*{前言}
\addcontentsline{toc}{chapter}{前言}

1996 年初, 我讲授了一门关于椭圆曲线 (elliptic curves) 的课程. 由于当时怀尔斯 (Wiles) 刚证明了费马大定理 (Fermat's Last Theorem) 不久, 且我承诺会解释他证明背后的一些思想, 这门课吸引了异常众多且背景各异的听众. 因此, 我尝试让所有仅具备标准研究生一年级课程知识的学生都能听懂这门课.

课程结束后, 我把课程期间发放的讲义整理成一个单独的文件, 做了少量修改, 并发布在互联网上, 获得了数万次的下载.

2006 年, 我对讲义进行了重写, 并以平装书的形式出版 (即本书的第 1 版). 对于这第 2 版, 我再次对全书进行了重写与更新, 同时保留了第 1 版的大部分编号.

除了其内在的数学趣味之外, 椭圆曲线的研究也是通往当代数论研究中一些更深层次领域的极佳切入点. 在本书中, 我试图将椭圆曲线理论置于这一更广阔的背景以及历史脉络之中.

在回顾椭圆曲线理论的发展历程时, 我深感其中有如此多的理论最初都源于具体的计算: Ramanujan 的计算启发了 Hecke 算子与Ramanujan猜想; Sato 的计算启发了广义形式的Sato--Tate猜想; Selmer 的计算启发了Cassels--Tate对偶定理; Birch 与 Stephens 的计算启发了 Gross--Zagier 公式; 当然, 还有Birch与 Swinnerton-Dyer 的计算, 启发了他们著名的 BSD 猜想及其推广.

谨在此向世界科技出版公司 (World Scientific Publishers) 的 Rochelle Kronzek 致谢, 感谢她鼓励我完成这第 2 版的编写工作.

\begin{flushright}
J.S. Milne
\end{flushright}